\documentclass[conference,onecolumn]{IEEEtran}
\IEEEoverridecommandlockouts
\usepackage{cite}
\usepackage{amsmath,amssymb,amsfonts}
\usepackage{algorithmic}
\usepackage{graphicx}
\usepackage{textcomp}
\usepackage{xcolor}
\usepackage{booktabs}
\usepackage{multirow}
\def\BibTeX{{\rm B\kern-.05em{\sc i\kern-.025em b}\kern-.08em
    T\kern-.1667em\lower.7ex\hbox{E}\kern-.125emX}}
\begin{document}

\title{Impact of Network Topology on Distributed Inference:\\Circulant Graphs as a Drop-In Replacement for 3D Torus Interconnects}

\author{
\IEEEauthorblockN{Nikhil Vemuri}
\IEEEauthorblockA{nvemuri@mit.edu}
\and
\IEEEauthorblockN{Brian Zhang}
\IEEEauthorblockA{placeholder@mit.edu}
}

\maketitle


\section{Project Overview}

This project is based on Example Project 5: Impact of Network Topology on Distributed Inference.

Modern accelerator clusters for deep learning inference distribute computation across tens to hundreds of chips connected by a high-speed interconnect. The topology of this interconnect, meaning which chips are directly connected to which, has a major impact on overall system performance. Google's TPU v4 pods, for example, use a 3D torus with 6 inter-chip interconnect (ICI) links per chip. But why a 3D torus? Is it actually the best topology for the communication patterns that arise in inference workloads, or is it simply convenient to manufacture?

This project studies that question by building a network topology model that plugs into AccelForge's mapping framework, evaluating several topologies across diverse workloads, and ultimately proposing an alternative topology, the circulant graph, that achieves 2.3x lower network latency than the 3D torus while using the same number of physical links per chip.

The key questions are:
\begin{enumerate}
    \item How much does the choice of network topology affect inference latency and energy, and where does the impact come from (which collective operations, which tensor transfers)?
    \item How sensitive is AccelForge's mapping to the network cost model, i.e., does feeding topology-aware costs back into the mapper change the spatial partitioning it selects?
    \item Can we find a topology that strictly dominates the 3D torus under the same hardware budget (same number of links per chip)?
\end{enumerate}

We hypothesized that the 3D torus would significantly outperform simpler topologies like meshes and rings due to its wraparound links reducing diameter and enabling ring AllReduce along each dimension independently. We also expected that higher-dimensional topologies might offer advantages by splitting data across more parallel rings. Both of these turned out to be correct, but the biggest surprise was the circulant graph, which achieves its advantage not by adding dimensions but by decomposing its links into edge-disjoint Hamiltonian cycles that can carry AllReduce traffic in parallel with zero contention.


\section{Technical Contributions}

\subsection{Network Topology Model}

We implemented a topology library (the \texttt{network\_topology} Python package) that models inter-chip communication for arbitrary graph topologies. Given an adjacency matrix, the model computes shortest-path routing, per-link byte loads for each collective operation (AllReduce, AllGather, ReduceScatter, Broadcast), and link-level congestion when multiple transfers run concurrently. The bottleneck link determines overall latency.

Each topology implements its own collective routing strategy. For torus topologies, AllReduce decomposes into independent ring AllReduces along each dimension, with per-link load $2(d-1)/d \cdot B$ for a ring of size $d$ carrying $B$ bytes. For the mesh, which lacks wraparound links, AllReduce uses a linear pipeline with per-link load $2B$ (a factor of $d/(d-1)$ worse than the ring). This distinction turned out to be critical: early versions of our code used the ring formula for both mesh and torus, which incorrectly showed them performing identically.

\subsection{Iterative Feedback Loop with AccelForge}

AccelForge's mapper optimizes spatial partitioning based on per-access energy and latency costs for each memory level. The ``NetworkMemory'' level in our architecture spec acts as a logical memory representing inter-chip communication. Initially, we set its costs to default values. After mapping, we extract the actual tensor transfers (which tensors are read/written from NetworkMemory, and how many bytes), infer the required collective operations from the mapping's spatial sharding decisions, evaluate them on the topology model, and feed the resulting per-byte energy and latency back as updated proxy costs. We repeat this until the proxy values converge (relative change below 5\%) or a maximum of 6 iterations.

This feedback loop is important because the mapper's partitioning decisions depend on the network cost, and the network cost depends on the partitioning. For most of our workloads, the loop converges in 2 to 3 iterations.

\subsection{Proposed Topology: Circulant Graph}

Our main contribution is proposing the circulant graph $C(64, \{1, 5, 17\})$ as a replacement for the 3D torus. In a circulant graph, chip $i$ connects to chips $(i \pm g) \bmod n$ for each generator $g$. With three generators and bidirectional links, each chip has exactly 6 links, matching the TPU v4 hardware budget.

The key property is that when each generator $g$ is coprime to $n$ (and here $\gcd(1,64)=1$, $\gcd(5,64)=1$, $\gcd(17,64)=1$), the cycle $0 \to g \to 2g \to \cdots \to (n{-}1)g \to 0 \pmod{n}$ visits every node exactly once, forming a Hamiltonian cycle. Since the three generators use disjoint edges, we get three edge-disjoint Hamiltonian cycles.

AllReduce then splits the data into three equal parts and runs a standard ring AllReduce on each cycle independently. Each ring has per-link load $2(n{-}1)/(3n) \cdot B$, compared to the torus's $2(d{-}1)/d \cdot B$ running sequentially across three dimensions of size $d{=}4$. The circulant's approach is better because (a) it parallelizes across all three rings simultaneously instead of running them in sequence, and (b) the rings never share links, so there is no congestion between them.

We also evaluated a 6D hypercube ($2^6 = 64$ chips), which has the same degree and diameter as the torus but achieves 1.5x faster AllReduce by decomposing into 6 rings of size 2. However, the circulant still beats it by a further 1.5x.

\subsection{Related Work}

Google's TPU v4 uses a 3D torus with reconfigurable ICI links \cite{tpuv4}. NVIDIA's NVSwitch provides all-to-all connectivity within a node but relies on fat-tree or dragonfly topologies across nodes \cite{nvswitch}. Circulant graphs have been studied in the interconnection network literature \cite{circulant_survey} as alternatives to hypercubes and tori, but to our knowledge, this is the first evaluation of their performance integrated with a full accelerator mapping framework.


\section{Evaluation}

We evaluate on a simulated system of 64 TPU v4 chips (275 TFLOPS BF16, 32 GiB HBM2 each, 6 ICI links at 45 GB/s per link, 4 pJ/bit/hop). AccelForge maps workloads to 8 chips, and we scale the network traffic to 64 chips. All topologies use the same chip architecture; only the interconnect changes.

\subsection{Experiment 1: Baseline Topology Comparison}

We compare five topologies (Circulant, 6D Hypercube, 3D Torus, 3D Mesh, Ring) across seven matrix multiplication workloads that span a range of shapes and sizes. The workloads are chosen so that the total working set exceeds 32 GiB, forcing AccelForge to shard tensors across chips and generating real inter-chip traffic.

\begin{table}[h]
\centering
\caption{Early results: network latency (ms) across topologies and workloads.}
\label{tab:latency}
\begin{tabular}{l r r r r r}
\toprule
Workload & Circ. & 6D Hyp. & Torus & Mesh & Ring \\
\midrule
Square 128K$^2$   & 8,017 & 12,217 & 18,325 & 24,434 & 24,052 \\
Wide 8K$\times$256K     & 2,004 & 3,054 & 4,581 & 6,108 & 6,013 \\
Tall 256K$\times$64K    & 16,035 & -- & 36,650 & 48,867 & 48,104 \\
VeryWide 2K$\times$256K & 251 & 382 & 573 & 764 & 752 \\
Rect 64K$\times$128K    & 8,017 & 12,217 & 18,325 & 24,434 & 24,052 \\
LargeSq 256K$^2$  & 32,069 & -- & 73,301 & 97,734 & 96,207 \\
VeryTall 256K$\times$8K & 3,150 & 6,204 & 7,731 & 15,557 & 9,163 \\
\bottomrule
\end{tabular}
\end{table}

The circulant wins every workload, with a consistent 2.29x speedup over the torus for AllReduce-dominated workloads. The 6D hypercube sits in between at 1.50x over torus. The mesh and ring perform similarly because, without wraparound links, the mesh's AllReduce degrades to linear pipelines that are only marginally better than the ring's.

\textbf{Why 2.29x?} The torus runs AllReduce sequentially across 3 dimensions of size 4, with per-link load $\frac{2 \cdot 3}{4} B = 1.5B$. The circulant splits data across 3 parallel rings of size 64, with per-link load $\frac{2 \cdot 63}{3 \cdot 64} B \approx 0.66B$, a ratio of $1.5 / 0.66 = 2.29$.

\subsection{Experiment 2: Energy Breakdown}

\begin{table}[h]
\centering
\caption{Early results: total energy in Joules (compute + network) for three representative workloads.}
\label{tab:energy}
\begin{tabular}{l r r r r r}
\toprule
Workload & Circ. & 6D Hyp. & Torus & Mesh & Ring \\
\midrule
Square 128K$^2$   & 5,480 & 10,019 & 8,330 & 13,397 & 5,480 \\
Wide 8K$\times$256K     & 1,398 & 2,532 & 2,110 & 3,377 & 1,398 \\
LargeSq 256K$^2$  & 34,717 & -- & 46,117 & 66,383 & 34,717 \\
\bottomrule
\end{tabular}
\end{table}

Compute energy is identical across topologies since the chip architecture and mapping are the same (the mapper converges to the same partitioning). The network energy differs because topologies with shorter average path lengths transmit fewer bit-hops per byte of data. The circulant's average hop count of 2.73 (vs. 3.05 for the torus) translates directly to lower energy per transfer.

An interesting finding is that the ring has the lowest energy despite having the worst latency. This is because ring AllReduce with 2 links moves data along short physical hops (always to the adjacent chip), minimizing bit-hop energy. The ring's problem is bandwidth, not energy. Network energy accounts for only 0.1 to 1\% of total energy for the ring, but 80 to 90\% of total latency.

\subsection{Experiment 3: GPT-3 Transformer Workloads}

Since AccelForge's mapper does not support multi-einsum workloads (e.g., full attention layers with softmax and residual connections), we separately evaluate GPT-3 transformer layers using direct topology modeling. We model tensor-parallel inference (Megatron-style) where each transformer layer requires two AllReduce operations (one after attention, one after FFN), each moving $\text{batch} \times \text{seq\_len} \times d_\text{model} \times 2$ bytes.

We evaluate GPT-3 175B (96 layers, $d$=12288), GPT-3 6.7B (32 layers, $d$=4096), and GPT-3 1.3B (24 layers, $d$=2048) across prefill and decode scenarios. In all cases, the circulant achieves 2.29x lower latency than the torus, matching the analytical prediction. The topology advantage is largest for long-context prefill (e.g., 32K tokens on GPT-3 175B generates 600 GB of AllReduce traffic per layer) and smallest for single-token decode, where the activation size is small and the fixed per-hop latency becomes relatively more significant.

\subsection{Experiment 4: Mapping Sensitivity}

The iterative feedback loop typically converges in 2 to 3 iterations with less than 5\% relative change in proxy values. For most workloads, the mapper selects the same spatial partitioning regardless of the topology-specific proxy costs, because the dominant factor in the mapping decision is the sheer size of the tensors relative to on-chip memory, not the network cost per byte. This means the topology comparison is largely apples-to-apples: the same partitioning, the same set of collective operations, and only the network cost model changes.

The exception is the VeryTall workload (256K$\times$8K), where the mapper generates three distinct transfers (broadcast of both activation tensors plus AllReduce of the output) instead of the single AllReduce seen in other workloads. This produces a 4.94x advantage for the circulant over the mesh, compared to the typical 3.05x, because broadcast routing is also more efficient on the circulant's low-diameter graph.


\section{Timeline and Work Distribution}

We have approximately three weeks remaining. The work so far and planned work is divided as follows.

\subsection{Completed Work}
\begin{itemize}
    \item \textbf{Nikhil:} Implemented the network topology library (Ring, Mesh, Torus3D, TorusND, CirculantHD), the AccelForge feedback loop integration, and the Slurm parallelization infrastructure on the MIT Engaging cluster. Ran the initial matmul sweep (7 workloads $\times$ 5 topologies) and the GPT-3 direct stress analysis.
    \item \textbf{Brian:} Implemented the AccelForge architecture configuration for the distributed TPU v4 system, the mapping serialization and test infrastructure, and contributed to the sweep driver script and workload selection.
\end{itemize}

\subsection{Remaining Work (Weeks 1--3)}
\begin{itemize}
    \item \textbf{Week 1 (Nikhil):} Run the full energy breakdown analysis with and without the network model to quantify the difference (Milestone 2 deliverable). Generate polished figures for all experiments.
    \item \textbf{Week 1 (Brian):} Investigate mapping sensitivity more deeply: compare the spatial partitionings chosen under different topology proxy costs and document cases where the mapping actually changes. Write up the evaluation methodology.
    \item \textbf{Week 2 (Nikhil):} Run scaling experiments (varying chip count from 16 to 256) to study how the circulant advantage changes with system size. Explore additional generator sets for the circulant to verify that $\{1, 5, 17\}$ is near-optimal.
    \item \textbf{Week 2 (Brian):} Expand the GPT-3 analysis to include data-parallel and pipeline-parallel strategies in addition to tensor parallelism. Compare the topology sensitivity across parallelism strategies.
    \item \textbf{Week 3 (Both):} Write the final report. Nikhil will cover the topology model, circulant proposal, and experimental results. Brian will cover the AccelForge integration, mapping sensitivity analysis, and the parallelism strategy comparison.
\end{itemize}

\begin{thebibliography}{00}
\bibitem{tpuv4} N. Jouppi et al., ``TPU v4: An Optically Reconfigurable Supercomputer for Machine Learning with Hardware Support for Embeddings,'' \textit{ISCA}, 2023.
\bibitem{nvswitch} NVIDIA, ``NVIDIA NVSwitch,'' NVIDIA Technical Brief, 2022.
\bibitem{circulant_survey} F. K. Hwang, ``A survey on multi-loop networks,'' \textit{Journal of Research of the National Bureau of Standards}, vol. 85, no. 4, 1980.
\end{thebibliography}

\end{document}
